\chapter*{Abstract}
Large Language Models produce fluent but factually incorrect responses up to 40\% of the time. This problem limits their usage in high-stakes situations. The current models cannot improve of its own because they operate statelessly without any knowledge accomulation. This research proposes CAEM (Confidence-Aware Episodic Memory with Self-Improvement), a system that progressively reduces hallucinations through verified rejection sampling and iterative fine-tuning. It integrates three core innovations: Verified episodic memory with multi-layer quality control, multi-signal confidence estimation for adaptive routing, and a closed feedback loop where verification creates training data that improves model generation. The verification module uses three different methods to check claims such as NLI entailment checking for factual claims, self-consistency for logical reasoning, and semantic entropy for confabulation detection. The goal is to achieve 85–90\% accuracy across all three methods. CAEM creates a virtuous improvement cycle through continuous cycles of iterative fine-tuning. In each cycle, verification produces high-quality training examples by fine-tuning. It allows to improve generation quality, and better generation produces quality training data. We apply three-tier routing for input problems, where exact matches are retrieved from memory, similar problems use the improved fine-tuned model, and novel problems use retrieval-augmented generation with multi layed verification. The study aims to reduce hallucinations by 20 to 30\% on established benchmarks (HotpotQA, StrategyQA, FEVER, TruthfulQA) without losing general model capabilities along with the goal to prove that even imperfect verification can provide monotonic improvement through data purity dynamics when baseline model accuracy exceeds random chance.


\vspace{0.5cm}

\noindent\textbf{Keywords:} Large Language Models, Hallucination Reduction, Episodic Memory, Confidence Estimation, Self-Improvement, Verification Systems, Continual Learning